\scanpage{56}
Si on a deux biicosaèdres orientés
$\{I_1,I'_1\}$ et $\{I_2,I'_2\}$, le choix d'une \deleted{icos. orient}
composante connexe du premier, et d'un
iso de celle-ci avec une comp. conn. de
l'autre, se prolonge de façon unique en
un isom des biicosaèdres orientés.

La paire \deleted{\unsure{$I$}} $(\widetilde I_\Phi,\widetilde I_{\Phi'})$ des rev. orientés
des deux struct. icosaèdrales \inserted{gauches} complémentaires,
est un biicos. orienté, et son groupe des
autom est celui du \deleted{groupe} biicosaèdre
gauche de départ.

Une autre façon de dire les choses est
celle-ci. Considérons un biicosaèdre droit
non orienté $I$. On lui associe l'icosa
complémentaire $I^{\unsure{\circ|o}}$, tel que pour $s,s'\in S=S(I)$
$\{s,s'\}$ soit arête de $I^\circ$ sss $\{s,\underline a(s')\}$ est arête de
$I$ (où $\underline a$ est l'antipodisme). On a donc
$I=I^{\circ\circ}$. \deleted{\unsure{On} le groupe des} Une structure
de biicosaèdre sur $S$ (\deleted{\unsure{sur|ensemble}} de douze sommets
est une paire de structures icosaédrales
$\{I,I'\}$ sur $S$, mutuellement complémentaires.

Le groupe des autom s'envoie dans
le groupe aux \inserted{(\unsure{sur}) le groupe des} permutations de $\{I,I'\}$, donc
\[
1\longrightarrow
\underset{\substack{\simeq\operatorname{Aut}(I')\\
\simeq\operatorname{Aut}^{+}(I)\times\pm1\\
\simeq\operatorname{Aut}^{+}(I')\times\pm1}}{\operatorname{Aut}(I)}
\longrightarrow\operatorname{Aut}(I,I')\longrightarrow\mathbb Z/2\mathbb Z
\longrightarrow\unsure{1|0}
\]
% archive sheet number: 55
